\documentclass[12pt]{article}
\usepackage{hyperref}

\usepackage{graphicx}

\usepackage{mathtools}
\DeclarePairedDelimiter\ceil{\lceil}{\rceil}
\DeclarePairedDelimiter\floor{\lfloor}{\rfloor}


\usepackage{amsmath}
\usepackage{amsfonts}
\usepackage{amssymb}

\usepackage{listings}
\lstset{language=C++,
                basicstyle=\ttfamily,
                keywordstyle=\color{blue}\ttfamily,
                stringstyle=\color{red}\ttfamily,
                commentstyle=\color{green}\ttfamily,
                morecomment=[l][\color{magenta}]{\#}
}

\usepackage{breqn}
\usepackage[]{algorithm2e}
\usepackage{physics}

\newcommand\fnurl[2]{%
  \href{#2}{#1}\footnote{\url{#2}}%
}

\newcommand{\Four}{\mathcal{F}}
\renewcommand{\(}{\left(}
\renewcommand{\)}{\right)}
%% \renewcommand{\{}{\left{}
%% \renewcommand{\}}{\right}}


\title{Real-to-Real FFT Notes for rocFFT}
\author{The rocFFT Team}



\def\no#1{\nobreak{#1}} % stop equation breaking in dmath

\begin{document}
\maketitle

\section{Introduction}

This document contains notes for computing real-to-real FFTs (like
sine transforms) for \texttt{rocFFT}.



\subsection{DST-I}


Given $\{x_n\}_{n=0}^{N-1}$,
\begin{dmath}
  X_k = \sum_{n=0}^{N-1}x_n \sin\left[\frac{\pi}{N+1}(n+1)(k+1)\right]
\end{dmath}
It would appear that this can be computed from a FFT of length $N$
with pre/post-processing.

This is equivalent to the negative imaginary part of the DFT of
\begin{dmath}
  0, x_0, x_1, \dots, x_{N-1}, 0, -x_{N-1}, \dots, -x_1, -x_0
\end{dmath}
and taking indices $1,\dots,N$.


\subsection{DST-IV}

Given $\{x_n\}_{n=0}^{N-1}$,
\begin{dmath}
  X_k = \sum_{n=0}^{N-1}x_n \sin\left[\frac{\pi}{N+1}\(n+\frac{1}{2}\)\(k+\frac{1}{2}\)\right]
\end{dmath}
It would appear that this can be computed from a FFT of length $N$
with pre/post-processing.

\end{document}
